% 中文译文：已发表靶点评估资源及本流程范围。
\begin{table}[!tbp]
\centering\small
\caption{Published target-assessment resources and the scope of this workflow.}\label{tab:framework-comparison}
\begin{tabularx}{\textwidth}{p{.23\textwidth}YY}\toprule
% 中文译文：资源；已发表范围；与本流程的连接
Resource & Published scope & Link to the present workflow \\ \midrule
% 中文译文：Open Targets：靶点—疾病证据整合与关联评分，提供异质关联和治疗证据。
Open Targets \citep{10.1093/nar/gkae1128} & Target--disease evidence integration and association scoring. & Provides heterogeneous association and therapeutic evidence. \\
% 中文译文：canSAR：癌症转化知识、化学和靶点可行性信息，为背景和可干预性评估提供参考。
canSAR \citep{10.1093/nar/gkaa1059} & Cancer translational knowledge, chemical and target-feasibility information. & Informs context and tractability assessment. \\
% 中文译文：Pharos/TCRD：人类靶点注释及知识成熟度评估，记录靶点注释和证据覆盖。
Pharos/TCRD \citep{10.1093/nar/gkac1033} & Human-target annotation and knowledge-maturity assessment. & Documents target annotation and evidence coverage. \\
% 中文译文：DepMap/Project Score：癌症模型功能依赖图谱，提供功能证据及来源特异标签构建。
DepMap/Project Score \citep{pacini_comprehensive_2024} & Cancer-model functional dependency mapping. & Provides functional evidence and source-specific label constructions. \\
% 中文译文：GOT-IT：靶点评估及验证规划建议，为背景特异后续研究提供指导。
GOT-IT \citep{emmerich_improving_2021} & Target-assessment and validation-planning recommendations. & Provides guidance for context-specific follow-up. \\
% 中文译文：本流程：双轴学习、来源约束及测量条件审核，将证据角色连接到排序、可能类别和带来源的证据卡。
Present workflow & Dual-axis learning, source restrictions and measurement-conditional review. & Executable linkage from evidence roles to rankings, possible classes and source-linked cards. \\
\bottomrule\end{tabularx}
% 中文译文：本表概括各资源已发表的功能及本流程实现的计算连接。
\parbox{\textwidth}{\footnotesize The table summarizes published functional scope and the computational linkage implemented in this workflow.}
\end{table}
